\documentclass[a4paper, 12pt]{article}

\usepackage[english, russian]{babel}
\usepackage[T2A]{fontenc}
\usepackage[utf8]{inputenc}
\setlength{\parindent}{1cm}
\usepackage{graphicx}
\usepackage{amsmath}
\usepackage{amssymb}
\usepackage{amsfonts}
\usepackage{array}
\usepackage{booktabs}
\usepackage{wrapfig}
\usepackage{caption} %заголовки плавающих объектов
\captionsetup[table]{justification=centering} %установки для заголовков таблиц
\captionsetup[figure]{justification=centering} %установки для заголовков таблиц


\begin{document}
	
	\begin{center}
		\subsubsection*{Расчётные задачи по дополнительным главам аналитической механики и электродинамики.}
		\subsubsection*{Выполнил студент третьего курса Физико-механического института, высшей школы фундаментальных физических исследований: Куликовских Илья Михайлович}
		\subsubsection*{группа: 5030302/30801}
		\subsubsection*{Преподаватель: Барсуков Дмитрий Петрович}
	\end{center}
	
	\quad \textbf{Задача Д1. } Найти уровни энергии для частицы массой \( m \), движущейся в одномерном потенциале
	
	\[
	U(x) = \sum_{n=1}^{N-1} \alpha \cdot \delta (x - na) \quad \text{при} \quad 0 \leq x \leq Na
	\]
	\[
	U(x) = +\infty \quad \text{при} \quad x < 0 \quad \text{и} \quad x > Na
	\]
	
	Где расположены зоны проводимости и запрещенные зоны?
	
	Чему равна эффективная масса частицы \( m_{\text{eff}} \) вблизи границы зоны проводимости?
	
	В пределе очень длинной решётки (\( N \to +\infty \)) найти, чему равна плотность числа состояний.
	
	\quad \textbf{Решение.} Необходимо найти волновые функции в конечном кристалле, для этого решим вспомогательную задачу на распределение зон в бесконечном кристалле, а потом перейдём к кристаллу конечных размеров. Напишем стационарное уравнение Шрёдингера: 
	
	\begin{equation}
		-\dfrac{\hbar^2}{2m} \dfrac{d^2 \psi}{dx^2} + \alpha\displaystyle\sum\limits_{n=-\infty}^{\infty} \delta(x-na) \psi = E\psi
	\end{equation}
	
	\quad Согласно теореме Блоха волновые функции в бесконечном кристалле имеют вид: 
	
	\begin{equation}
		\psi(x) = e^{ikx} u_k(x) \quad u_k (x) = u_k (x+a)
	\end{equation}
	
	\quad где $k$ - волновое число, $\hbar k$ - квазиимпульс, $a$ - вектор трансляции решётки. 
	
	\quad В пространстве между дельта функциями волновая функция равна волновой функции свободной частица. Для учёта решётки введём специальные граничные условия - условия сшивания на дельта-потенциале и условия периодичности из теоремы Блоха. Имеем:
	
	\begin{center}
		$\psi = C_1 e^{iqa} + C_2 e^{-iqa} \quad q = \sqrt{\dfrac{2mE}{\hbar^2}}$ 
		
		Г. у. : $\begin{cases}
			\psi_k (na-0) = e^{ika} \psi((n-1)a+0) \\
			\\
			\dfrac{d\psi_k}{dx} \bigg|_{x\to na-0} = e^{ika}\dfrac{d\psi}{dx}\bigg|_{x\to (n-1)a+0} \\
			\\
			\dfrac{d\psi_k }{dx} \bigg|_{x\to na+0} - \dfrac{d\psi_k}{dx}\bigg|_{x \to na-0} = \dfrac{2m\alpha}{\hbar^2} \psi_k (na) \\ 
		\end{cases}$
		
	\end{center}
	
	\quad Объединяя условия на скачок производной и на периодичность производной из теоремы Блоха получаем два уравнения из которых можно построить систему для определения собственных чисел: 
	
	\begin{equation}
		\begin{cases}
			\dfrac{d\psi_k}{dx} \bigg|_{x=(n-1)a} - e^{-ika} \dfrac{d\psi_k}{dx}\bigg|_{x=na} = \dfrac{2m\alpha}{\hbar^2} \psi_k((n-1)a) \\
			\psi_k ((n-1)a) = e^{ika} \psi_k(na)  
		\end{cases}
	\end{equation}
	
	\begin{equation}
		\begin{cases}
			C_1 (e^{iqa} - e^{ika}) + C_2 (e^{-iqa} - e^{ika} ) = 0 \\
			iq C_1 \left(e^{i(q-k)a} - 1 - \dfrac{2m\alpha}{iq\hbar^2}\right) + iqC_2 \left(1-e^{-i(q+k)a} - \dfrac{2m\alpha}{iq\hbar^2}\right) = 0
		\end{cases}
	\end{equation}
	
	\begin{center}
		$(e^{iqa} - e^{ika} ) \left(1 - e^{i(q+k)a} + \dfrac{2m\alpha}{iq\hbar^2}\right) - \left(e^{i(q-k)a} -1 + \dfrac{2m\alpha}{iq\hbar^2} \right) (e^{-iqa} - e^{ika}) = 0$
		
		$2(e^{iqa} + e^{-iqa}) - 2(e^{ika} + e^{-ika}) + \dfrac{2m\alpha}{iq\hbar^2} (e^{iqa} - e^{-iqa}) = 0$
	\end{center}
	
	\begin{equation}
		\cos(qa)  + \dfrac{m\alpha}{q\hbar^2} \sin(qa) = \cos(ka) 
	\end{equation}
	
	\quad Уравнением (5) описываются разрешённые и запрещённые значения $q$, а как следствие и значения $E$ т.e. 
	
	\begin{center}
		$ -1 < \cos(qa)  + \dfrac{m\alpha}{q\hbar^2} \sin(qa) < 1$ 
	\end{center}
	
	\quad Решая это уравнение численно можно найти распределение зон в бесконечном кристалле. Теперь перейдём к кристаллу конечных размеров. Для этого изменим немного граничные условия - добавим однородные условия Дирихле на границе кристалла: 
	
	\begin{equation}
		\text{Г. у.: } \begin{cases}
			\psi(0) = \psi(Na) = 0 \\ 
			\psi(na+0) = \psi(na-0) \\
			\dfrac{d\psi }{dx} \bigg|_{x = na+0} - \dfrac{d\psi }{dx} \bigg|_{x = na-0} = \dfrac{2 m \alpha}{\hbar^2}\psi(na)
		\end{cases}
	\end{equation}
	
	\begin{center}
		$\psi = A_n e^{iq(x - (n-1)a)} + B_n e^{-iq(x - (n-1)a)} \quad x \in ((n-1)a; na)$
	\end{center}
	
	\quad Подставляем в эти граничные условия и получаем систему линейных уравнений: 
	
	\begin{equation}
		\begin{cases}
			A_n e^{iqa} + B_n e^{-iqa} = A_{n+1} + B_{n+1} \\
			A_{n+1} - B_{n+1} - A_n e^{iqa} \left(1 - \dfrac{2m\alpha}{iq\hbar^2} \right) + B_n e^{-iqa} \left(1 + \dfrac{2m\alpha}{iq\hbar^2} \right) = 0 \\
			A_1 + B_1 = 0 \\
			A_N e^{iqa} + B_N e^{-iqa} = 0
		\end{cases}
	\end{equation}
	
	\quad Если рассмотреть только первые два уравнения этой системы, то можно вывести рекуррентное соотношение между коэффициентами $A_n$, $B_n$: 
	
	\begin{equation}
		\begin{cases}
			A_n e^{iqa} + B_n e^{-iqa} = A_{n+1} + B_{n+1} \\
			A_{n+1} - B_{n+1} - A_n e^{iqa} \left(1 - \dfrac{2m\alpha}{iq\hbar^2} \right) + B_n e^{-iqa} \left(1 + \dfrac{2m\alpha}{iq\hbar^2} \right) = 0 
		\end{cases}
	\end{equation}
	
	\begin{equation}
		\begin{cases}
			A_{n+1} = A_n e^{iqa} \left(1 - \dfrac{m\alpha}{iq\hbar^2}\right) + B_n e^{-iqa} \dfrac{m\alpha}{iq\hbar^2} \\
			B_{n+1} = -A_n e^{iqa} \dfrac{m\alpha}{iq\hbar^2} + B_n e^{-iqa} \left(1+ \dfrac{m\alpha}{iq\hbar^2}\right) 
		\end{cases} \implies
	\end{equation}
	
	\begin{equation}
		\implies \text{т.е. } \left(\begin{array}{c}
			A_{n+1} \\
			B_{n+1} 
		\end{array} \right)= \underbrace{\left(\begin{array}{cc}
				e^{iqa} \left(1 - \frac{m\alpha}{iq\hbar^2}\right) & e^{-iqa} \frac{m\alpha}{iq\hbar^2} \\
				-e^{iqa} \frac{m\alpha}{iq\hbar^2} & e^{-iqa} \left(1+ \frac{m\alpha}{iq\hbar^2}\right) 
			\end{array}\right) }_{T} \cdot \left(\begin{array}{c}
			A_n \\
			B_n
		\end{array}\right)
	\end{equation}
	
	\quad где $T$ - матрица трансляции. 
	
	\quad В силу того, что матрица трансляции связывает два соседних коэффициента при её возведении в степень $(N-1)$ можно задать явную связь между первыми и последними коэффициентами $A_n$, $B_n$. 
	
	\begin{equation}
		\left(\begin{array}{c}
			A_N \\
			B_N 
		\end{array}\right) = T^{N-1} \left(\begin{array}{c}
			A_1 \\
			B_1
		\end{array}\right)
	\end{equation}
	
	\quad Учитывая уравнения из условия (7) и систему (11) мы можем однозначно определить все коэффициенты, а следовательно и ВФ в кристалле. Для вычисления этих коэффициентов в явном виде диагонализуем матрицу $T$: 
	
	\begin{center}
		$\det (T - \lambda E ) = 0 \implies \left|\begin{array}{cc}
			e^{iqa} \left(1 - \frac{m\alpha}{iq\hbar^2}\right) - \lambda& e^{-iqa} \frac{m\alpha}{iq\hbar^2} \\
			-e^{iqa} \frac{m\alpha}{iq\hbar^2} & e^{-iqa} \left(1+ \frac{m\alpha}{iq\hbar^2}\right) -\lambda
		\end{array}\right| = 0$
	\end{center} 
	
	\begin{equation}
		\lambda^2 - 2\lambda\left(\dfrac{m\alpha}{q\hbar^2} \sin(qa) + \cos(qa) \right) + 1 = 0
	\end{equation}
	
	\quad Анализируя формулу (12) можно понять, что коэффициенты $A_n$, $B_n$ зависят от собственных чисел матрицы $T$. 
	
	\quad 1) $D = 4\left(\dfrac{m\alpha}{q\hbar^2} \sin(qa) + \cos(qa) \right)^2 - 4 > 0 $
	
	\quad Нетрудно заметить, что в этом случае собственные числа $T$ вещественны и $E$ соответствует запрещённой зоне (по следствию из формулы (5)). 
	
	2) $D = 4\left(\dfrac{m\alpha}{q\hbar^2} \sin(qa) + \cos(qa) \right)^2 - 4 < 0$
	
	\quad Тут собственные числа комплексны и энергия соответствует разрешённым зонам, т.е. 
	
	\begin{equation}
		\dfrac{m\alpha}{q\hbar^2} \sin(qa) + \cos(qa)  < \cos(ka)
	\end{equation}
	
	\quad Найдём модуль собственного числа $\lambda$ как модуль комплексного числа: 
	
	\begin{center}
		$\lambda_{1, 2} = \dfrac{2\left(\dfrac{m\alpha}{q\hbar^2} \sin(qa) + \cos(qa)\right) \pm 2i \sqrt{1 -\left(\dfrac{m\alpha}{q\hbar^2} \sin(qa) + \cos(qa)\right)^2 }}{2}$
		
		$\lambda_{1,2} = \left(\dfrac{m\alpha}{q\hbar^2} \sin(qa) + \cos(qa)\right) \pm i \sqrt{1 - \left(\dfrac{m\alpha}{q\hbar^2} \sin(qa) + \cos(qa)\right)^2}$
	\end{center}
	
	\begin{equation}
		|\lambda| = \left(\dfrac{m\alpha}{q\hbar^2} \sin(qa) + \cos(qa)\right)^2 + 1 - \left(\dfrac{m\alpha}{q\hbar^2} \sin(qa) + \cos(qa)\right)^2 = 1
	\end{equation}
	
	\begin{equation}
		\implies \lambda = e^{\pm i\varphi} \quad |\lambda| = 1
	\end{equation}
	
	\quad Рассмотрим теперь собственные векторы матрицы трансляции $T$. Положим, что собственный вектор это $\left( A_1  \quad B_1 \right)^{T}$:
	
	\begin{center}
		$T\left(\begin{array}{c}
			A_1 \\
			B_1
		\end{array}\right) = \left(\begin{array}{c}
			A_2 \\
			B_2
		\end{array}\right) = \lambda\left(\begin{array}{c}
			A_1 \\
			B_1 
		\end{array}\right)$
	\end{center} 
	
	\quad Таким образом:
	
	\begin{equation}
		\left(\begin{array}{c}
			A_n \\
			B_n
		\end{array}\right) = \lambda^{n-1}\left(\begin{array}{c}
			A_1 \\
			B_1 
		\end{array}\right)
	\end{equation}
	
	\quad Из теоремы Блоха: $\psi(x+na) = e^{ikna} \psi(x) \Leftrightarrow \psi_{n+1} (x) = e^{ikna} \psi_1 (x) $.
	Волновые функции в кристалле таким образом: $\lambda_{1,2} = e^{\pm ika} \quad \varphi = ka$
	
	\begin{equation}
		\psi(x) = \begin{cases}
			A_1 e^{iqx} + B_1 e^{-iqx} \quad \forall x \in (0, a)\\
			e^{ika} \left(A_1 e^{iqx} + B_1 e^{-iqx} \right) \quad \forall x \in (a, 2a) \\
			\dots\\
			e^{ik(N-1)a} \left(A_1 e^{iqx} + B_1 e^{-iqx} \right) \quad \forall x \in ((N-1)a, Na)  
		\end{cases}
	\end{equation}
	
	\begin{center}
		$\begin{cases}
			A_1 + B_1 = 0 \\
			A_N e^{iqa} + B_N e^{-iqa} = 0
		\end{cases} \implies e^{ikNa} - e^{-ikNa} = 0$
	\end{center}
	
	\begin{equation}
		\sin(kNa) = 0 \implies kNa = \pi n \quad k_n = \dfrac{\pi n }{Na} \quad n \in \mathbb{Z}
	\end{equation}
	
	\quad Формула (18) выражает условие квантования квазиимпульса, а формула (13) устанавливает связь между $k$ и $q$, т.е. решая трансцендентное уравнение (13) при заданных $k$ получаем уровни энергии $E$. Применим теперь условие нормировки и получим: 
	
	\begin{center}
		$A_1 + B_1 = 0 \implies \psi_1(x) = 2iA_1 \sin(qa)$ 
	\end{center}
	
	\begin{equation}
		\displaystyle\int\limits_{0}^{a}  |\psi_1(x)|^2 dx = 4A_1^2 \displaystyle\int\limits_{0}^{a} \sin^2 (qa) dx = 2A_1^2 - A_1^2 \dfrac{2qa - \sin(qa)}{q} 
	\end{equation}
	
	\begin{equation}
		1 = N\cdot A_1^2 \cdot \dfrac{2qa - \sin(2qa)}{q} \implies A_1 = \dfrac{1}{\sqrt{N}} \cdot \sqrt{\dfrac{q}{2qa - \sin(2qa)}}
	\end{equation}
	
	\quad Таким образом волновые функции кристалла полностью определены через матрицу $T$ и начальный коэффициент $A_1$. 
	
	\quad Для нахождения эффективной массы $m_{\text{eff}}$ используем формулу:
	
	\begin{equation}
		m_{\text{eff}} = \hbar^2 \left( \dfrac{d^2E}{dk^2} \right)^{-1}
	\end{equation}
	
	\quad Запишем дисперсионное соотношение (5) в виде:
	
	\[
	\cos(ka) = \cos(qa) + \dfrac{m\alpha}{\hbar^2 q} \sin(qa), \quad q = \sqrt{\dfrac{2mE}{\hbar^2}}
	\]
	
	\quad Дифференцируем обе части по $k$:
	
	\[
	-a\sin(ka) = \left[ -a\sin(qa) + \dfrac{m\alpha}{\hbar^2} \left( \dfrac{\cos(qa)}{q} - \dfrac{\sin(qa)}{q^2} \right) \right] \dfrac{dq}{dk}
	\]
	
	\quad Учитывая, что $\dfrac{dE}{dk} = \dfrac{\hbar^2 q}{m} \dfrac{dq}{dk}$, получаем:
	
	\[
	\dfrac{dE}{dk} = \dfrac{\hbar^2 \sin(ka)}{a} \left[ \sin(qa) - \dfrac{m\alpha}{\hbar^2 q} \left( \dfrac{\cos(qa)}{q} - \dfrac{\sin(qa)}{q^2} \right) \right]^{-1}
	\]
	
	\quad Вторую производную находим как:
	
	\[
	\dfrac{d^2E}{dk^2} = \dfrac{d}{dk} \left( \dfrac{dE}{dk} \right)
	\]
	
	\quad Вблизи границы зоны ($k=0$ или $k=\pi/a$) производные упрощаются. Для $k=0$ (дно зоны) в приближении слабого потенциала:
	
	\[
	m_{\text{eff}} \approx m \left( 1 + \dfrac{2m\alpha a}{\hbar^2} \right)
	\]
	
	\quad Для потолка зоны ($k=\pi/a$) эффективная масса отрицательна:
	
	\[
	m_{\text{eff}} \approx -m \left( 1 - \dfrac{2m\alpha a}{\hbar^2} \right)
	\]
	
	\quad Таким образом, эффективная масса зависит от квазиимпульса и знака потенциала $\alpha$.
	
	\quad Число состояний в интервале $dk$:
	
	\[
	dN = \dfrac{Na}{2\pi} dk
	\]
	
	\quad Плотность состояний по энергии:
	
	\[
	\rho(E) = \dfrac{dN}{dE} = \dfrac{Na}{2\pi} \dfrac{dk}{dE}
	\]
	
	\quad Используя связь $\dfrac{dk}{dE} = \left( \dfrac{dE}{dk} \right)^{-1}$ и выражение для $\dfrac{dE}{dk}$ из предыдущего раздела, получаем:
	
	\[
	\rho(E) = \dfrac{Na}{2\pi\hbar} \left| \dfrac{dk}{dE} \right| = \dfrac{Na}{2\pi\hbar} \left| \sin(qa) - \dfrac{m\alpha}{\hbar^2 q} \left( \dfrac{\cos(qa)}{q} - \dfrac{\sin(qa)}{q^2} \right) \right|^{-1}
	\]
	
	\quad Эта формула определяет плотность состояний в бесконечном кристалле.
	
	
	\quad \textbf{Задача Д2.} Найти коэффициент прохождения \( D \) для частицы массой \( m \), движущейся в одномерном потенциале  
	
	\[
	U(x) = \sum_{n=1}^{N} \alpha \cdot \delta (x - na)
	\]
	
	При каких значениях энергии мы будем иметь \( D = 1 \)?
	
	\quad \textbf{Решение.} На основе полного решения задачи Д1 можно сразу написать волновые функции кристалла. Слева от потенциалов волновая функция является суперпозицией падающей волны и отражённой волны: 
	
	\begin{equation}
		\psi(x) = e^{iqx} + R e^{-iqx} \quad \forall x < a
	\end{equation}
	
	\begin{center}
		где $R^2$ - коэффициент отражения 
	\end{center}
	
	\quad Далее волновые функции повторяют ВФ из выражения для предыдущей задачи. 
	
	\begin{equation}
		\psi (x) = A_n e^{iq(x-(n-1)a)} + B_n e^{-iq(x-(n-1)a)} \quad \forall x \in [(n-1)a, na]
	\end{equation}
	
	\quad И напоследок волновые функции справа от кристалла составляют только лишь прошедшие кристалл волны:
	
	\begin{equation}
		\psi (x) = D e^{iqx} \quad \forall x > Na 
	\end{equation}
	
	\begin{center}
		где $D^2$ - коэффицент прохождения.
	\end{center}
	
	\quad Рассмотрим теперь первый дельта потенциал: 
	
	\begin{equation}
		\begin{cases}
			\psi_1 (x)\bigg|_{x=a-} = \psi_2 (x) \bigg|_{x=a+} \\
			\dfrac{d\psi }{dx} \bigg|_{x = na+0} - \dfrac{d\psi }{dx} \bigg|_{x = na-0} = \dfrac{2 m \alpha}{\hbar^2}\psi(na)
		\end{cases} \implies
	\end{equation}
	
	\begin{equation}
		\implies \begin{cases}
			A_2 + B_2 = e^{iqa} + Re^{-iqa} \\
			\left(1 - \dfrac{2m\alpha}{iq\hbar^2}\right) A_2 - \left(1 + \dfrac{2m\alpha}{iq\hbar^2}\right) B_2 = e^{iqa} - Re^{-iqa}
		\end{cases}
	\end{equation}
	
	\begin{center}
		$\left(\begin{array}{cc}
			1 & 1 \\
			\left(1 - \frac{2m\alpha}{iq\hbar^2}\right) & -\left(1 + \frac{2m\alpha}{iq\hbar^2}\right) 
		\end{array}\right) \left(\begin{array}{c}
			A_2 \\
			B_2
		\end{array} \right) = \left(\begin{array}{cc}
			e^{iqa} & e^{-iqa} \\
			e^{iqa} & -e^{-iqa} 
		\end{array}\right) \left(\begin{array}{c}
			1 \\
			R
		\end{array}\right)$
		
		
		$\left(\begin{array}{c}
			1 \\
			R
		\end{array}\right) =  \left(\begin{array}{cc}
			e^{iqa} & e^{-iqa} \\
			e^{iqa} & -e^{-iqa} 
		\end{array}\right)^{-1}\left(\begin{array}{cc}
			1 & 1 \\
			\left(1 - \frac{2m\alpha}{iq\hbar^2}\right) & -\left(1 + \frac{2m\alpha}{iq\hbar^2}\right) 
		\end{array}\right) \left(\begin{array}{c}
			A_2 \\
			B_2
		\end{array} \right) $
		
		Пусть $P = \left(\begin{array}{cc}
			e^{iqa} & e^{-iqa} \\
			e^{iqa} & -e^{-iqa} 
		\end{array}\right)$, тогда 
		
		$(P|E) = \left(\begin{array}{cc|cc}
			e^{iqa} & e^{-iqa} & 1 & 0 \\
			e^{iqa} & -e^{-iqa} & 0 & 1
		\end{array}\right) = \left(\begin{array}{cc|cc}
			1 & 0 & \frac{1}{2}e^{-iqa} & \frac{1}{2}e^{-iqa} \\
			0 & 1 & \frac{1}{2}e^{iqa} & -\frac{1}{2}e^{iqa}
		\end{array}\right) = (E|P^{-1})$
		
		\begin{multline*}
			\left(\begin{array}{cc}
				\frac{1}{2}e^{-iqa} & \frac{1}{2}e^{-iqa} \\
				\frac{1}{2}e^{iqa} & -\frac{1}{2}e^{iqa}
			\end{array}\right)
			\left(\begin{array}{cc}
				1 & 1 \\
				\left(1 - \frac{2m\alpha}{iq\hbar^2}\right) & -\left(1 + \frac{2m\alpha}{iq\hbar^2}\right) 
			\end{array}\right) = \\
			= \left(\begin{array}{cc}
				e^{-iqa} \left(1 - \frac{m\alpha}{2iq\hbar^2}\right) & -e^{-iqa}\frac{m\alpha}{2iq\hbar^2} \\
				e^{iqa} \frac{m\alpha}{2iq\hbar^2} & e^{iqa}\left(1+\frac{m\alpha}{2iq\hbar^2} \right)
			\end{array}\right)
		\end{multline*}
	\end{center}
	
	\begin{equation}
		\left(\begin{array}{c}
			1 \\
			R
		\end{array}\right) = \underbrace{\left(\begin{array}{cc}
		e^{-iqa} \left(1 - \frac{m\alpha}{2iq\hbar^2}\right) & -e^{-iqa}\frac{m\alpha}{2iq\hbar^2} \\
		e^{iqa} \frac{m\alpha}{2iq\hbar^2} & e^{iqa}\left(1+\frac{m\alpha}{2iq\hbar^2} \right)
		\end{array}\right) }_{T_L} \left(\begin{array}{c}
			A_2 \\
			B_2 
		\end{array}\right)
	\end{equation}
	
	\quad Тем самым мы получили связь между коэфицентом отражения и коэфицентами в первой зоне кристалла. Теперь проделаем аналогичные действия с самым правым дельта потенциалом: 

	\begin{center}
		$\psi_N (x) = A_N e^{iq(x-(N-1)a)} + B_N e^{-iq(x - (N-1)a)} \quad \forall x \in [(N-1)a, Na] $
		
		$\psi_{N+1} (x) = D e^{iqx} + 0\cdot e^{-iqx} \quad \forall x > Na$
		
		$\dfrac{d\psi_{N+1}}{dx} = iqDe^{iq(x - (N-1)a)} - iq\cdot 0 \cdot e^{-iq(x-(N-1)a)}$
		
		$\begin{cases}
			\psi_N (x) \bigg|_{Na+} = \psi_{N+1} (x) \bigg|_{Na-} \\ 
			\dfrac{d\psi_{N+1}}{dx}\bigg|_{Na+} - \dfrac{d\psi_{N}}{dx}\bigg|_{Na-} = \dfrac{2m\alpha}{\hbar^2} \psi(Na)
		\end{cases}$
		
		$\begin{cases}
			A_N e^{iqa} + B_N e^{iqa} = De^{iqNa} \\
			De^{iqNa} - A_N e^{iqa} + B_N e^{-iqa} = \dfrac{2m\alpha}{iq\hbar^2} (A_N e^{iqa} + B_N e^{-iqa}) 
		\end{cases}$
		
		
		$\begin{cases} 
			A_N e^{iqa} + B_N e^{-iqa} = D e^{iqNa} + 0\cdot e^{-iqNa} \\
			A_N e^{iqa} \left(1+ \dfrac{2m\alpha}{iq\hbar^2} \right) - B_N e^{-iqa} \left(1 - \dfrac{2m\alpha}{iq\hbar^2}\right) = De^{iqNa} - 0 \cdot e^{-iqNa} 
		\end{cases}$
		
		\begin{multline*}
			\left(\begin{array}{c}
				D \\
				0
			\end{array}\right) = \left(\begin{array}{cc}
				e^{iqNa} & e^{-iqNa} \\
				e^{iqNa} & -e^{-iqNa} 
			\end{array}\right)^{-1} \cdot \\
			\cdot \left(\begin{array}{cc}
				e^{iqa} & e^{-iqa} \\
				e^{iqa} \left(1 + \frac{2m\alpha}{iq\hbar^2}\right) & -e^{-iqa}\left(1 - \frac{2m\alpha}{iq\hbar^2}\right) 
			\end{array}\right) \left(\begin{array}{c}
				A_N \\
				B_N
			\end{array}\right)
		\end{multline*}
	\end{center}
	\begin{equation}
		\left(\begin{array}{c}
			D \\
			0
		\end{array}\right) =  \underbrace{\left(\begin{array}{cc}
			e^{-iqNa + iqa} \left(1 + \frac{m\alpha}{iq\hbar^2}\right) & e^{-iqNa - iqa} \cdot \frac{m\alpha}{iq\hbar^2} \\
			- e^{iqNa + iqa} \cdot \frac{m\alpha}{iq\hbar^2} & e^{iqNa - iqa} \left(1 - \frac{m\alpha}{iq\hbar^2}\right)
		\end{array}\right)}_{T_R} \left(\begin{array}{c}
			A_N \\
			B_N \\
		\end{array}\right)
	\end{equation}
	
	\begin{center}
		$\left(\begin{array}{c}
			A_N \\
			B_N
		\end{array}\right) = T^{N-2} \left(\begin{array}{c}
			A_2 \\
			B_2 
		\end{array}\right) $
	\end{center}
	
	\begin{equation}
		\implies \left(\begin{array}{c}
			D \\
			0
		\end{array}\right) = T_R T^{N-2} \left(\begin{array}{c}
			A_2 \\
			B_2 
		\end{array}\right) = T_R T^{N-2} T_L^\dag \left(\begin{array}{c}
		1 \\
		R 
		\end{array}\right)
	\end{equation}
	
	\quad Система линейных уравнений (29) полностью и однозначно связывает коэфиценты прохождения и отражения, а также позволяет определить их через другие параметры этой системы. Для опредления явного вида $D$ достаточно будет возвести матрицу $T$ в степень $N-2$.  В силу унитарности любых преобразований координат получается: 
	
	\begin{equation}
		T^{N-2} = \dfrac{\sin\left((N-2)ka\right)}{\sin{(ka)}} T - \dfrac{\sin\left((N-3)ka\right)}{\sin{(ka)}} E
	\end{equation}
	
	\quad Пропуская дальнейшее упрощение: 
	
	\begin{center}
		\begin{multline*}
			\left(\begin{array}{cc}
				\cos[(N-1)ka] + i\displaystyle\frac{\beta}{2}\frac{\sin[(N-2)ka]}{\sin ka} & i\displaystyle\frac{\beta}{2}e^{-2iqa}\frac{\sin[(N-2) ka]}{\sin ka} \\[10pt]
				-i\displaystyle\frac{\beta}{2}e^{2iqa}\frac{\sin[(N-2) ka]}{\sin ka} & \cos[(N-1) ka] - i\displaystyle\frac{\beta}{2}\frac{\sin[(N-2) ka]}{\sin ka}
			\end{array}\right) \cdot \\ \cdot
			\left(\begin{array}{c}
				1 \\ R
			\end{array}\right)
			=
			\left(\begin{array}{c}
				D \\ 0
			\end{array}\right)
		\end{multline*}
	\end{center}
	
	\quad Выражает $D$ из системы линейный уравнений выходит: 
	
	\begin{equation}
		D = \frac{1}{\cos[(N-1) ka] - i\displaystyle\frac{\beta}{2}\frac{\sin[(N-2) ka]}{\sin ka}}
	\end{equation}
	
	\quad Коэффицент прохождения $D^2$ равен в таком случае: 
	
	\begin{equation}
		|D|^2 = \frac{1}{1 + \left(\displaystyle\frac{\beta}{2}\right)^2 \displaystyle\frac{\sin^2[(N-2) ka]}{\sin^2 ka}}
	\end{equation}
	
	\quad Если взять $D^2 = 1 \implies  \left(\displaystyle\frac{\beta}{2}\right)^2 \displaystyle\frac{\sin^2[(N-2) ka]}{\sin^2 ka} = 0$
	
	\begin{equation}
		\begin{cases}
			\sin^2 [(N-2)ka] = 0 \\
			\sin^2 ka = 0 
		\end{cases} \Leftrightarrow \begin{cases}
			(N-2)ka = \pi n \\
			ka \neq \pi s
		\end{cases} \quad \forall n, s \in \mathbb{Z}
	\end{equation}
	
	\quad Соотношением (33) описываются всевозможные значения $k$, $q$ при которых коэффицент прохождения мог бы быть равен 1. 
	
	

	
\end{document}